\begin{afterword}
\summaryhead

The post opens with a plucked chicken offered as a human ``by definition''. People who argue by definition, the author says, start from visible features, find them in a dictionary definition, and draw their conclusion. But features both sides accept are rarely what the dispute is about. The legitimate move from visible to hidden properties, from Socrates's human shape to his vulnerability to hemlock, is a probabilistic inference from clusters of similar things, and no definition changes those clusters.

If someone reports seeing Socrates with herbologists preparing an antidote, the argument must turn to biochemistry. Insisting that humans are mortal ``by definition'' then refuses to use evidence already seen, like calling a coin 50\% heads after looking at it. People add the phrase, the post says, when the default inference has been doubted, or when X does not resemble the central members of Y, as in ``atheism is, by definition, a religion''. The author's tally over thirteen years suggests that ``by definition'' outside mathematics is among the worst signs of a flawed argument, so readers should drop the phrase.

\responsehead

Most of the post is sound. The point that the premise of an argument by definition is usually uncontested, and so cannot be what is at stake, is a real diagnosis. So is the account of the legitimate inference: from observed features to hidden ones through the cluster of similar things, with new evidence allowed to override it when it bears on the prediction. The coin and Frodo analogies state the rule exactly, and the post limits the rule itself: nine fingers do not change the hemlock prediction.

The weak part is the claim about when people use the phrase. The post says people add ``by definition'' ``on exactly those occasions'' when the default inference is in doubt, or when a connotation is to be sneaked in. The cases are the author's own. The herbologists are a hypothetical, and the atheism argument is the author's wording of the opposing case, with no one quoted. The evidence offered is a tally kept for thirteen years, which the post itself describes as eyeballing, ``not with literal statistics.''

The advice has a scope the post leaves implicit. It fits everyday argument, where a category is used to predict. Where a definition is itself a rule that assigns consequences, as in law, arguing from it is legitimate; an American appeals court classed atheism as a religion for First Amendment purposes in 2005.

In short: a sound account of how evidence overrides a category, with a claim about when people say ``by definition'' that rests on the author's informal tally and on hypothetical cases.
\end{afterword}
